\chapter{LaTeX Usage Examples}
\label{app:latex}
\exampleparagraph{This appendix is a quick reference for writing the report.
Remove its input line from \texttt{main.tex} before submission if it is not needed.}

\section{Citations and cross-references}
Add entries to \texttt{references.bib}, then cite their keys:
\begin{lstlisting}[language=TeX,numbers=none]
\citet{wolfLam1991locality} describe a locality algorithm.
Related work considers loop transformations
\citep{wolfLam1991locality,mckinleyCarrTseng1996}.
\citep[pp.~30--31]{wolfLam1991locality}

\cref{chap:background}
\cref{fig:access-pattern}
\cref{tab:results}
\cref{lst:baseline}
\cref{eq:row-major}
\Cref{fig:access-pattern} % Capitalise at the start of a sentence.
\cref{lst:baseline,lst:interchanged} % Multiple references.
\end{lstlisting}
Put a figure or table label immediately after its caption. Put equation labels
inside the equation environment. Compile the document and bibliography using
the provided build command so that citations and references are resolved.
Use \verb|\cref{...}| within a sentence and \verb|\Cref{...}| at its beginning.
The object name (such as ``figure'', ``equation'', or ``listing'') is inserted
automatically, so do not write it again before the command. Bibliographic
citations still use \verb|\citet| and \verb|\citep|.

\section{Mathematics}
\begin{lstlisting}[language=TeX,numbers=none]
Inline notation: $\vect{i}\in\Z^n$.
\begin{equation}
  W(B) = 3eB^2 \leq \alpha C
  \label{eq:your-new-label}
\end{equation}
\begin{align}
  x &= a+b, \\
  y &= c+d.
\end{align}
\end{lstlisting}
The template loads \texttt{amsmath}, \texttt{amssymb}, \texttt{mathtools}, and
\texttt{amsthm}; the background chapter also demonstrates matrices and a definition.

% Start the second guide page here so the complete example stays together.
\clearpage
\section{Code and external source files}
For short snippets, use an inline listing environment:
\begin{latexexample}
\begin{lstlisting}[language=Python,caption={Your caption.},
                   label={lst:your-label}]
for tile in tiles:
    process(tile)
\end{lstlisting}
\end{latexexample}
For reusable code, keep it in \texttt{code/} and include it as follows:
\begin{lstlisting}[language=TeX,numbers=none]
\lstinputlisting[language=C,caption={Your caption.},
                 label={lst:your-file}]{code/your-file.c}
\end{lstlisting}
\texttt{listings} needs no external syntax-highlighting tool or shell escape.
Give a referenced listing both a caption and a label so it has its own number.
The example C files use C99 variable-length array parameters and are intended
as explanations of loop order rather than a benchmark harness.

\section{Figures and tables}
Save images under \texttt{figures/}. Prefer PDF for vector plots and PNG or JPEG
for raster images. A basic figure looks like this:
\begin{lstlisting}[language=TeX,numbers=none]
\begin{figure}[htbp]
  \centering
  \includegraphics[width=0.8\linewidth]{your-figure.pdf}
  \caption{Explain what the figure shows and cite its source.}
  \label{fig:your-label}
\end{figure}
\end{lstlisting}
% \Cref{chap:example} shows side-by-side panels with \texttt{subcaption}.
% Use \texttt{booktabs} rules for tables and \texttt{tabularx} for wrapping columns;
% use \texttt{longtable} if a table must span pages. Let figures and tables float
% with \texttt{[htbp]} unless their exact position is essential.

\section{Small formatting details}
Escape reserved characters: \verb|\%|, \verb|\&|, \verb|\_|, and \verb|\#|.
Use \verb|~| for a non-breaking space, as in \verb|Section~1|.
Smart references such as \verb|\cref{fig:access-pattern}| supply the name and spacing.
Use \verb|\footnote{...}| for footnotes\footnote{This is an example footnote.},
\verb|\url{https://example.org}| for URLs, and
\verb|\href{URL}{text}| for links.
Use \verb|\todo{...}| for visible draft reminders.
